\newpage
\chapter{PENDAHULUAN} \label{Bab I}

\section{Latar Belakang Masalah} \label{I.Latar Belakang}
Kemajuan teknologi informasi dan komunikasi telah mengubah cara masyarakat memproduksi, mendistribusikan, dan mengonsumsi informasi digital. Dalam konteks tersebut, media visual seperti citra dan video menempati posisi penting karena sering dipersepsikan sebagai representasi langsung dari realitas. Namun, perkembangan kecerdasan buatan generatif juga meningkatkan risiko manipulasi informasi visual dalam skala besar. \textit{Global Risks Report 2026} yang diterbitkan \textit{World Economic Forum} menempatkan \textit{misinformation and disinformation} pada peringkat kedua risiko global dalam horizon dua tahun, menunjukkan bahwa gangguan terhadap integritas informasi merupakan isu strategis yang semakin mendesak \cite{WorldEconomicForumWEF2026}. Salah satu bentuk manipulasi visual yang paling menonjol adalah \textit{deepfake}, yaitu media sintetis berbasis \textit{deep learning} yang mampu menghasilkan konten wajah dan video yang sangat realistis, sehingga berpotensi digunakan untuk penipuan, manipulasi opini, dan penyalahgunaan identitas digital \cite{Tolosana2020,BhaskardeepKhaund2025}. \par

Ancaman \textit{deepfake} tidak lagi terbatas pada ranah hiburan atau eksperimen teknis, tetapi telah berkembang menjadi persoalan keamanan informasi dan kepercayaan digital. Tolosana dkk. menjelaskan bahwa kemajuan teknik manipulasi wajah, khususnya yang ditopang \textit{deep learning} dan GAN, telah menghasilkan konten palsu yang sangat meyakinkan serta menimbulkan tantangan serius bagi sistem deteksi \cite{Tolosana2020}. Sejalan dengan itu, Khaund menunjukkan bahwa media sintetis modern dapat melemahkan keandalan mekanisme autentikasi yang bergantung pada biometrik dan verifikasi dokumen, sehingga \textit{deepfake} menjadi ancaman yang relevan bagi ekosistem identitas digital \cite{BhaskardeepKhaund2025}. Kondisi ini menegaskan perlunya pengembangan sistem deteksi otomatis yang tidak hanya akurat pada lingkungan terkontrol, tetapi juga tangguh ketika dihadapkan pada variasi kualitas media di dunia nyata \cite{BhaskardeepKhaund2025}. \par

Dalam perkembangannya, pendekatan deteksi \textit{deepfake} dapat dibedakan menjadi pendekatan spasial dan pendekatan spasiotemporal. Pendekatan spasial, yang secara operasional bekerja pada tingkat \textit{frame} tunggal, berfokus pada ekstraksi artefak visual statis seperti ketidakwajaran tekstur, batas wajah, dan inkonsistensi detail lokal \cite{Tolosana2020,Cao2022}. Pendekatan ini relatif efisien dan telah menunjukkan performa yang baik pada banyak dataset terkontrol. Akan tetapi, semakin realistisnya generator \textit{deepfake} membuat artefak visual statis menjadi semakin halus, sehingga detektor yang hanya bergantung pada bukti spasial berisiko mengalami penurunan kemampuan generalisasi \cite{Tolosana2020,Cao2022}. Dengan kata lain, peningkatan kualitas generasi media sintetis menuntut strategi deteksi yang tidak hanya mengandalkan petunjuk visual pada satu bingkai secara terpisah. \par

Sebagai alternatif, pendekatan spasiotemporal memanfaatkan hubungan antarbingkai untuk mendeteksi inkonsistensi gerakan atau koherensi temporal yang tidak alami pada video \textit{deepfake}. Gu dkk. menunjukkan bahwa deteksi video \textit{deepfake} pada praktiknya berkembang dalam dua jalur besar, yaitu solusi berbasis \textit{frame} dan solusi berbasis video, serta menekankan pentingnya pemodelan \textit{spatial-temporal inconsistency} dalam video hasil manipulasi \cite{Gu2021}. Demikian pula, Haliassos dkk. menunjukkan bahwa representasi gerakan mulut yang bersifat spasiotemporal dapat meningkatkan kemampuan generalisasi dan ketahanan terhadap berbagai distorsi umum \cite{Haliassos2021}. Meskipun demikian, keunggulan pendekatan spasiotemporal tidak selalu bersifat mutlak, karena isyarat temporal juga dapat terdegradasi oleh kompresi, \textit{post-processing}, dan kualitas video yang rendah, yang lazim dijumpai pada distribusi konten daring \cite{Haliassos2021}. \par

Tantangan tersebut menjadi semakin penting ketika sistem deteksi diuji pada skenario \textit{in-the-wild}. Zi dkk. memperkenalkan WildDeepfake sebagai dataset dunia nyata yang terdiri atas 7.314 sekuens wajah yang diekstraksi dari 707 video \textit{deepfake} yang dikumpulkan sepenuhnya dari internet \cite{Zi2024}. Mereka juga menunjukkan bahwa WildDeepfake merupakan dataset yang lebih menantang dibandingkan dataset yang lebih terkontrol, karena performa sejumlah detektor dasar dapat menurun secara drastis pada dataset ini \cite{Zi2024}. Temuan tersebut menunjukkan bahwa evaluasi pada data yang terkontrol belum cukup untuk merepresentasikan kondisi operasional sebenarnya. Oleh karena itu, pertanyaan mengenai pendekatan mana yang lebih tangguh, apakah pendekatan spasial atau pendekatan spasiotemporal, masih relevan untuk diuji secara empiris pada data \textit{deepfake} dunia nyata \cite{Gu2021,Haliassos2021,Zi2024}. \par

Berdasarkan uraian tersebut, penelitian ini berfokus pada analisis komparatif klasifikasi video \textit{deepfake} dengan pendekatan spasial dan spasiotemporal menggunakan dua model dari keluarga Transformer, yaitu Vision Transformer (ViT) dan encoder VideoMAE. ViT merepresentasikan pendekatan Transformer untuk citra dengan memproses urutan \textit{image patches} secara langsung \cite{dosovitskiy2021imageworth16x16words}, sedangkan VideoMAE merupakan \textit{video transformer} yang dirancang untuk mempelajari representasi video secara efisien \cite{tong2022videomaemaskedautoencodersdataefficient}. Dengan membandingkan kedua pendekatan tersebut pada dataset WildDeepfake, penelitian ini bertujuan untuk mengevaluasi secara objektif kontribusi relatif informasi spasial dan temporal terhadap ketahanan klasifikasi video \textit{deepfake} pada kondisi dunia nyata. Hasil penelitian diharapkan dapat memberikan dasar empiris yang lebih kuat mengenai pendekatan yang lebih sesuai untuk pengembangan sistem deteksi \textit{deepfake} yang \textit{robust} dan relevan secara praktis. \par

\section{Rumusan Masalah} \label{I.Rumusan Masalah}
\begin{enumerate}[noitemsep]
	\item Bagaimana merancang dan mengimplementasikan model klasifikasi video \textit{deepfake} menggunakan pendekatan spasial berbasis Vision Transformer (ViT) dan pendekatan spasiotemporal berbasis encoder VideoMAE pada dataset WildDeepfake?
	\item Bagaimana mengukur serta membandingkan performa pendekatan spasial berbasis Vision Transformer (ViT) dan pendekatan spasiotemporal berbasis encoder VideoMAE berdasarkan \textit{accuracy}, \textit{precision}, \textit{recall}, \textit{F1-score}, AUC-ROC, dan \textit{confusion matrix}?
	\item Pendekatan manakah di antara model spasial berbasis Vision Transformer (ViT) dan model spasiotemporal berbasis encoder VideoMAE yang paling efektif untuk klasifikasi video \textit{deepfake} pada skenario \textit{in-the-wild}?
\end{enumerate}

\section{Tujuan Penelitian} \label{I.Tujuan}
Berdasarkan rumusan masalah yang telah diuraikan, tujuan penelitian ini adalah sebagai berikut. \par

\begin{enumerate}[noitemsep]
	\item Merancang dan mengimplementasikan model klasifikasi video \textit{deepfake} menggunakan pendekatan spasial berbasis Vision Transformer (ViT) dan pendekatan spasiotemporal berbasis encoder VideoMAE pada dataset WildDeepfake.
	\item Mengukur serta membandingkan performa pendekatan spasial berbasis Vision Transformer (ViT) dan pendekatan spasiotemporal berbasis encoder VideoMAE berdasarkan \textit{accuracy}, \textit{precision}, \textit{recall}, \textit{F1-score}, AUC-ROC, dan \textit{confusion matrix}.
	\item Menentukan pendekatan yang paling efektif di antara model spasial berbasis Vision Transformer (ViT) dan model spasiotemporal berbasis encoder VideoMAE untuk klasifikasi video \textit{deepfake} pada skenario \textit{in-the-wild}.
\end{enumerate}

\section{Batasan Masalah} \label{I.Batasan}
Batasan yang dimaksud disini ialah batasan dari penelitian tugas akhir yang dilakukan. Batasan masalah ditujukan agar tugas akhir yang dilakukan tidak terlalu luas, dan menjadi lebih realistis untuk diselesaikan. \par

\section{Manfaat Penelitian} \label{I.Manfaat}
Manfaat tugas akhir yang dilakukan didefinisikan sebagai manfaat yang diperoleh ketika tugas akhir telah selesai dilakukan. Manfaat dapat berupa manfaat untuk: \textbf{mahasiswa, program studi teknik informatika, ITERA, masyarakat, dunia akademisi, dan dunia penelitian}. \par


\section{Sistematika Penulisan} \label{I.Sistematika}
Sistematika penulisan berisi pembahasan apa yang akan ditulis disetiap Bab. Sistematika pada umumnya berupa paragraf yang setiap paragraf mencerminkan bahasan setiap Bab. \par

\subsection*{Bab I}
Bab ini berisikan penjelasan latar belakang dari topik penelitian yang berlangsung, rumusan masalah dari masalah yang dihadapi pada penjelasan di latar belakang, tujuan dari penelitian, batasan dari penelitian, manfaat dari hasil penelitian, dan sistematika penulisan tugas akhir. \par

\subsection*{Bab II}
Bab ini membahas mengenai tinjauan pustaka dari penelitan terdahulu dan dasar teori yang berkaitan dengan penelitian ini.

\subsection*{Bab III}
Bab ini berisikan penjelasan alur kerja sistem, alat dan data yang digunakan, metode yang digunakan, dan rancangan pengujian.

\subsection*{Bab IV}
Bab ini membahas hasil implementasi dan pengujian dari penelitian yang dilakukan, serta analisis dan evaluasi yang dapat dipetik dari hasil.

\subsection*{Bab V}
Bab ini membahas kesimpulan dari hasil penelitian dan juga saran untuk penelitian selanjutnya.
